\documentclass{letter}
%\usepackage{fancyhdr}

% \oddsidemargin=1in
% \evensidemargin=1in
% \textwidth=7.0in
%\topmargin=-2.2in
% \textheight=9in

\usepackage{geometry}
\geometry{margin=1in}

% \pagestyle{fancy}
% \fancyhf{}
% \rhead{Ryan Giordano}
% \lhead{Research Statement}
% \rfoot{Page \thepage}

\usepackage{etoolbox}

\usepackage{microtype}
\usepackage{graphicx}
%\usepackage{subfigure}
\usepackage{booktabs} % for professional tables
\usepackage{xcolor}
\usepackage[hidelinks=True]{hyperref}
\usepackage{xargs}[2008/03/08]

% Documentation
% http://ftp.math.purdue.edu/mirrors/ctan.org/macros/latex/contrib/refstyle/refstyle.pdf
\usepackage{refstyle}
\usepackage{varioref} % Use refstyle instead of varioref directly.

\usepackage{amsmath}
\usepackage{amssymb}
\usepackage{amsfonts}
\usepackage{amsthm}
\usepackage{mathrsfs} % For mathscr
\usepackage{mathtools}

%\usepackage[authoryear]{natbib}
%\bibliographystyle{apalike}


\begin{document}
\begin{letter}{2021 Alvarez-Hopper Selection Committee\\
Lawrence Berkeley National Laboratory\\
Berkeley, CA, 94720}


\opening{Dear Alvarez-Hopper Selection Committee,}

I am writing to enthusiastically apply for the 2021 Alvarez postdoctoral
fellowship.  Through my work as an engineer at Google, my PhD in statistics at
Berkeley, and my current postdoctoral position at the Massachusetts Institute of
Technology, my academic ambition has consistently been to bring conceptually
simple, theoretically sound, and computationally tractable tools to bear on the
large scale data science problems of the twenty-first century.  The US
Department of Energy (DOE), and the Berkeley Lab in particular, are currently
aggressively pursuing new frontiers in data science, as expressed, for example,
by the Advanced Scientific Computing Advisory Committee Subcommittee on AI/ML
report of September 2020.  The Lab's expressed need for scalable and reliable
uncertainty quantification and model interrogation are a perfect fit for my
research, and I am confident that close collaboration between Berkeley
Lab scientists and myself would produce both valuable science and
methodological advances in data science.

My work on large-scale data problems, both at Google and on the Celeste project
(producing astronomical catalogs on the NERSC supercomputing cluster), have
acquainted me with the difficulty of applying traditional statistical ideas at
large scale.  My research has been devoted to easing these difficulties. My PhD
began with an effort to overcome the fact that classical Bayesian Markov Chain
Monte Carlo (MCMC) techniques would not scale to the problems I encountered at
Google. This research question led me to develop the linear response correction
for variational Bayes approximations, a perturbative correction which provides
reliable, approximate Bayesian inference on problems far too massive for MCMC
(like the construction of astronomical catalogs). During the second half of my
PhD, I found that similar perturbative ideas provide scalable solutions to a
wide array of practical machine learning problems involving uncertainty
quantification or model checking, including cross validation, prior sensitivity,
adversarial data sensitivity, and sensitivity to prior or model specification.

Good ideas are made better by being put into practice, and I believe that the my
research will be best developed and advanced in the context of massive,
practically relevant scientific applications.  In addition to revisiting the
production of astronomical catalogs in light of my variational Bayes work, both
the High Energy Physics and Nuclear Physics groups of Berkeley Lab have
expressed explicit need for uncertainty quantification as part of the ``AI for
Science Initiative'' in the above-mentioned report.  As the report states,
``While  AI algorithms have proliferated,  general  techniques  for
understanding  their accuracy  and stability in complex environments have not.''
I believe that collaboration between statisticians such as myself and scientists
is the best way to fill this need, and it is for this reason that I am applying
to become an Alvarez fellow.


\closing{Sincerely,}

\vspace{-5em}
\includegraphics[width=2.5in]{sig.png}

Ryan Giordano
% Department of Computer Science\\
% Massachusetts Institute of Technology\\
% Cambridge, MA\\
% (805) 501-6754 (cell)\\
% rgiordan@mit.edu\\
% \url{https://rgiordan.github.io/}
\end{letter}

% \bibliography{references}
% \bibliographystyle{plainnat}

\end{document}
